\section{Objetivos}
\label{sec:objetivos}

\subsection{Objetivo geral}
Compreender e implementar, na prática, as operações matemáticas necessárias para transformar coordenadas bidimensionais entre o Sistema de Coordenadas do Mundo, o Sistema de Coordenadas Normalizadas do Dispositivo (NDC) e as Coordenadas do Dispositivo, utilizando como representação visual um pixel de cor verde na tela.

\subsection{Objetivos específicos}
\begin{itemize}
    \item \textbf{Modelagem de dados:} Estruturar as classes necessárias para representar pontos cartesianos e os limites das janelas de visualização.
    \item \textbf{Normalização:} Desenvolver os algoritmos de conversão entre o Sistema de Coordenadas do Mundo e o Sistema de Coordenadas Normalizadas do Dispositivo (NDC), considerando os intervalos $[0,1]$ e $[-1,1]$.
    \item \textbf{Mapeamento para a tela:} Implementar a conversão das coordenadas NDC para as Coordenadas do Dispositivo, transformando as posições normalizadas em coordenadas de pixels de acordo com as dimensões da tela.
    \item \textbf{Aplicação interativa:} Construir uma interface gráfica utilizando a biblioteca Java Swing, permitindo ao usuário visualizar e validar o funcionamento das transformações implementadas.
\end{itemize}
